% ============================== PROBLEM 3 ============================
\settab{PROBLEM 3}{4}
\renewcommand{\chapribbontext}{Problem 3 \;·\; 5 pts}
\chapter[Problem 3: Radial Functions]{Radial Functions}\label{ch:p3}

\begin{examnote}[title={The mock}]
On $\Omega=\R^2\setminus\{(0,0)\}$, let $r=\sqrt{x^2+y^2}$ and $u(x,y)=\varphi(r)$. Harmonic means
$\Delta u=u_{xx}+u_{yy}=0$.
\end{examnote}

\begin{idea}
A \textbf{radial} function only depends on the distance $r$ to the origin (ripples around a stone).
Part A turns $u_{xx}+u_{yy}$ into $\varphi''+\frac1r\varphi'$; Part B solves it. After Q1 every answer is 2 or 3 lines.
\end{idea}

\begin{remember}[title={The only two things to memorise}]
\[
  r_x=\frac xr,\quad r_y=\frac yr \qquad\qquad \Delta u=\varphi''(r)+\frac1r\,\varphi'(r)
\]
\end{remember}

% -------------------------------------------------------------------
\section{Q1. First and second derivatives}



\begin{recipe}
$u$ depends on $x$ only through $r$: derivative of $\varphi$ times derivative of $r$.
\begin{enumerate}
  \item $u_x=\varphi'(r)\cdot r_x=\varphi'(r)\frac xr$.
  \item Side line: $\frac{\partial}{\partial x}\Big(\frac xr\Big)=\frac{r-x\cdot\frac xr}{r^2}=\frac{r^2-x^2}{r^3}=\frac{y^2}{r^3}$.
  \item Product rule on $\varphi'(r)\cdot\frac xr$: \ (derivative of first)$\cdot$second $+$ first$\cdot$(derivative of second).
\end{enumerate}
\end{recipe}

\begin{writethis}
$r_x=\frac xr$, $r_y=\frac yr$. \ Chain rule: \ $u_x=\varphi'(r)\dfrac xr$, \ $u_y=\varphi'(r)\dfrac yr$.

Since $\frac{\partial}{\partial x}\big(\frac xr\big)=\frac{y^2}{r^3}$ and $\frac{\partial}{\partial x}\big[\varphi'(r)\big]=\varphi''(r)\frac xr$:
\[
  u_{xx}=\varphi''(r)\frac xr\cdot\frac xr+\varphi'(r)\frac{y^2}{r^3}=\varphi''(r)\frac{x^2}{r^2}+\varphi'(r)\frac{y^2}{r^3}.
\]
Swapping $x$ and $y$: \ $u_{yy}=\varphi''(r)\dfrac{y^2}{r^2}+\varphi'(r)\dfrac{x^2}{r^3}$.
\end{writethis}

\begin{watch}\watchout The derivative of $\varphi'(r)$ in $x$ is $\varphi''(r)\cdot\frac xr$, not just $\varphi''(r)$.\end{watch}

% -------------------------------------------------------------------
\section{Q2. The Laplacian}

\begin{writethis}
Add, and use $x^2+y^2=r^2$: \
$\Delta u=\varphi''\dfrac{x^2+y^2}{r^2}+\varphi'\dfrac{x^2+y^2}{r^3}=\varphi''(r)+\dfrac1r\varphi'(r).$
\end{writethis}

% -------------------------------------------------------------------
\section{Q3. The logarithm}

\begin{writethis}
$v=\ln r$, so $\varphi(r)=\ln r$, \ $\varphi'=\frac1r$, \ $\varphi''=-\frac1{r^2}$. \
$\Delta v=-\frac1{r^2}+\frac1r\cdot\frac1r=0$: harmonic on $\Omega$.

Limit: as $(x,y)\to(0,0)$, $r\to0^+$ and $\ln r\to-\infty$. So $\lim v=-\infty$ (no finite limit).
\end{writethis}

% -------------------------------------------------------------------
\section{Q4 and Q5. Find every radial harmonic function}

\begin{idea}
Multiply $\varphi''+\frac1r\varphi'=0$ by $r$: you get $(r\varphi')'=0$ (product rule). Zero derivative means constant. Twice: $A$, then $B$.
\end{idea}

\begin{writethis}
\textbf{Q4.} $u$ harmonic, so $\varphi''(r)+\frac1r\varphi'(r)=0$ for every $r>0$. Then
\[
  \big(r\varphi'(r)\big)'=\varphi'(r)+r\varphi''(r)=r\Big(\varphi''+\tfrac1r\varphi'\Big)=0 .
\]
\seeproof{8} Zero derivative on the interval $(0,+\infty)$, so $r\varphi'(r)=A$, a constant.

\textbf{Q5.} $\big(\varphi(r)-A\ln r\big)'=\varphi'(r)-\frac Ar=\frac{r\varphi'(r)-A}{r}=0$, so
$\varphi(r)-A\ln r=B$. Hence
\[
  \boxed{\,u(x,y)=A\ln\sqrt{x^2+y^2}+B\,}.
\]
\end{writethis}

% -------------------------------------------------------------------
\section{Q6. When is the limit at the origin finite?}

\noindent\begin{minipage}[c]{0.47\linewidth}
\begin{writethis}
As $(x,y)\to(0,0)$, $r\to0^+$ and $\ln r\to-\infty$, so $u=A\ln r+B$ tends to: \
$-\infty$ if $A>0$; \ $+\infty$ if $A<0$; \ $B$ if $A=0$.
The limit is finite \textbf{if and only if $A=0$} (any $B$, not $A=B=0$), and then it equals $B$.
\end{writethis}
\end{minipage}\hfill
\begin{minipage}[c]{0.5\linewidth}\centering

\begin{tikzpicture}
\begin{axis}[bp, width=6.6cm, height=4.8cm, xmin=0, xmax=3.3, ymin=-2.6, ymax=4.2,
   xlabel=$r$, ylabel={$u$}, xtick={1,2,3}, ytick={-2,-1,1,2,3}]
  \addplot[tealD, line width=1.4pt, domain=0.04:3.2, samples=120] {ln(x)+1};
  \addplot[coralD, line width=1.4pt, domain=0.04:3.2, samples=120] {-ln(x)+1};
  \addplot[tealDD, line width=1.4pt, dashed, domain=0:3.2] {1};
  \node[font=\scriptsize, text=tealD, anchor=west] at (axis cs:2.2,2.05) {$A=1$: $-\infty$ at $0$};
  \node[font=\scriptsize, text=coralD, anchor=west] at (axis cs:1.9,-0.45) {$A=-1$: $+\infty$ at $0$};
  \node[font=\scriptsize, text=tealDD, anchor=south west] at (axis cs:2.2,1.02) {$A=0$};
\end{axis}
\end{tikzpicture}
\figcap{$u=A\ln r+B$ with $B=1$. Only $A=0$ arrives at the centre with a finite value.}

\end{minipage}

